% Slide type: plain
\begin{frame}{Section 5.4.1 ─ チャネル間隔の計算式}
\begin{itemize}
  \item 5.4.1: 隣接キャリア間の中心周波数間隔（公称値）
  \item 公称式: $\mathrm{Spacing} = (\mathrm{BW}_1 + \mathrm{BW}_2)/2$ ─ 端がちょうど接する距離
  \item 実配置: 公称値をチャネルラスター $\Delta F_{\mathrm{Raster}}$ の倍数に丸め
  \item NR-LTE 隣接時: NR/LTE 両ラスターを同時に満たす位置に丸め
  \item オペレーターは公称値より小さい間隔も選択可能（ガードバンド犠牲で帯域効率優先）
\end{itemize}
\end{frame}
